\chapter{Software Project Estimation}

\section{Function Point Analysis (FPA) Methodology}

Software project estimation is the quantitative discipline in software engineering that predicts the size, effort, duration, and staffing required to build a software system. Function Point Analysis (FPA), established by Allan Albrecht and standardized by the International Function Point Users Group (IFPUG), is an industry-standard indirect metric that measures the functional size of software from the user's business viewpoint, independent of programming language, development framework, or deployment platform.

Function Point Analysis categorizes system functionalities into five standard functional components:

\begin{enumerate}
  \item \textbf{External Inputs (EI):} Elementary processes that process data or control information entering the system from external entities, updating internal logical files.
  \item \textbf{External Outputs (EO):} Elementary processes that present derived, processed, or calculated data to external entities, typically involving mathematical transformation or report generation.
  \item \textbf{External Inquiries (EQ):} Elementary processes that retrieve data from internal logical files and present it to external entities without modifying system state or performing derived mathematical calculations.
  \item \textbf{Internal Logical Files (ILF):} User-identifiable groups of logically related data maintained completely within the system boundary through system business processes.
  \item \textbf{External Interface Files (EIF):} User-identifiable groups of logically related data referenced by the system but maintained entirely within external applications.
\end{enumerate}

\section{External Inputs (EI) Identification}

\subsection{Table 6.1: External Inputs (EI) Complexity Classification}

\begin{table}[htbp]
\centering
\small
\resizebox{\textwidth}{!}{%
\begin{tabular}{l l l l l l}
\toprule
\textbf{EI ID} & \textbf{Functional Process Description} & \textbf{Data Element Types (DET)} & \textbf{File Types Referenced (FTR)} & \textbf{Assessed Complexity} & \textbf{Weight} \\
\midrule
\textbf{EI-1} & Customer Registration & 12 & 1 (CUSTOMER) & Average & 4 \\
\textbf{EI-2} & Customer Login / Credential Authentication & 4 & 1 (CUSTOMER) & Low & 3 \\
\textbf{EI-3} & Customer Profile Update & 8 & 1 (CUSTOMER) & Low & 3 \\
\textbf{EI-4} & Loan Application Submission & 10 & 2 (CUSTOMER, LOAN\_APPLICATION) & Average & 4 \\
\textbf{EI-5} & Supporting Document Upload & 6 & 2 (LOAN\_APPLICATION, DOCUMENT) & Low & 3 \\
\textbf{EI-6} & Bank Employee Underwriting Review / Decision & 6 & 2 (LOAN\_APPLICATION, BANK\_EMPLOYEE) & Average & 4 \\
\textbf{EI-7} & Document Verification Status Update & 5 & 2 (DOCUMENT, BANK\_EMPLOYEE) & Low & 3 \\
\textbf{EI-8} & Loan Sanction and Disbursal Authorization & 8 & 2 (LOAN\_APPLICATION, LOAN) & Average & 4 \\
\textbf{EI-9} & Loan Repayment / EMI Transaction Submission & 8 & 2 (LOAN, REPAYMENT) & Average & 4 \\
\textbf{EI-10} & Admin Create / Update Loan Product Category & 7 & 1 (LOAN\_CATEGORY) & Low & 3 \\
\textbf{EI-11} & Admin Employee Account Management & 9 & 1 (BANK\_EMPLOYEE) & Average & 4 \\
\textbf{EI-12} & Customer Password Reset Request & 4 & 1 (CUSTOMER) & Low & 3 \\
\textbf{Total EI Points} & - & - & - & - & \textbf{42} \\
\bottomrule
\end{tabular}%
}
\end{table}

\section{External Outputs (EO) Identification}

\subsection{Table 6.2: External Outputs (EO) Complexity Classification}

\begin{table}[htbp]
\centering
\small
\resizebox{\textwidth}{!}{%
\begin{tabular}{l l l l l l}
\toprule
\textbf{EO ID} & \textbf{Functional Process Description} & \textbf{Data Element Types (DET)} & \textbf{File Types Referenced (FTR)} & \textbf{Assessed Complexity} & \textbf{Weight} \\
\midrule
\textbf{EO-1} & Formal Loan Sanction Letter Generation (PDF) & 16 & 3 (CUSTOMER, APP, LOAN) & High & 7 \\
\textbf{EO-2} & Complete EMI Amortization Schedule Calculation & 12 & 2 (LOAN, REPAYMENT) & High & 7 \\
\textbf{EO-3} & Monthly Repayment Receipt Generation (PDF) & 10 & 2 (CUSTOMER, REPAYMENT) & Average & 5 \\
\textbf{EO-4} & Loan Rejection Formal Notification Letter & 8 & 2 (CUSTOMER, APP) & Average & 5 \\
\textbf{EO-5} & Executive Loan Portfolio Analytics Summary & 18 & 4 (APP, LOAN, REPAY, CUST) & High & 7 \\
\textbf{EO-6} & Delinquency / Overdue Repayments Audit Report & 14 & 3 (CUSTOMER, LOAN, REPAY) & High & 7 \\
\textbf{EO-7} & Final Loan Closure Certificate (NOC) & 10 & 2 (CUSTOMER, LOAN) & Average & 5 \\
\textbf{Total EO Points} & - & - & - & - & \textbf{43} \\
\bottomrule
\end{tabular}%
}
\end{table}

\section{External Inquiries (EQ) Identification}

\subsection{Table 6.3: External Inquiries (EQ) Complexity Classification}

\begin{table}[htbp]
\centering
\small
\resizebox{\textwidth}{!}{%
\begin{tabular}{l l l l l l}
\toprule
\textbf{EQ ID} & \textbf{Functional Process Description} & \textbf{Data Element Types (DET)} & \textbf{File Types Referenced (FTR)} & \textbf{Assessed Complexity} & \textbf{Weight} \\
\midrule
\textbf{EQ-1} & Customer Application Status Inquiry & 6 & 2 (CUSTOMER, LOAN\_APPLICATION) & Low & 3 \\
\textbf{EQ-2} & Customer Active Loan Balance \& History Inquiry & 8 & 2 (LOAN, REPAYMENT) & Average & 4 \\
\textbf{EQ-3} & Employee Pending Applications Queue Search & 9 & 2 (CUSTOMER, LOAN\_APPLICATION) & Average & 4 \\
\textbf{EQ-4} & Employee Document Inspection View & 6 & 1 (DOCUMENT) & Low & 3 \\
\textbf{EQ-5} & Loan Schemes Catalog \& Rate Inquiry & 7 & 1 (LOAN\_CATEGORY) & Low & 3 \\
\textbf{EQ-6} & Admin User Audit Search Inquiry & 8 & 2 (CUSTOMER, BANK\_EMPLOYEE) & Average & 4 \\
\textbf{Total EQ Points} & - & - & - & - & \textbf{21} \\
\bottomrule
\end{tabular}%
}
\end{table}

\section{Internal Logical Files (ILF) Identification}

\subsection{Table 6.4: Internal Logical Files (ILF) Complexity Classification}

\begin{table}[htbp]
\centering
\small
\resizebox{\textwidth}{!}{%
\begin{tabular}{l l l l l l}
\toprule
\textbf{ILF ID} & \textbf{Logical File / Table Name} & \textbf{Record Element Types (RET)} & \textbf{Data Element Types (DET)} & \textbf{Assessed Complexity} & \textbf{Weight} \\
\midrule
\textbf{ILF-1} & \texttt{CUSTOMER} & 1 & 13 & Low & 7 \\
\textbf{ILF-2} & \texttt{LOAN\_APPLICATION} & 1 & 11 & Low & 7 \\
\textbf{ILF-3} & \texttt{BANK\_EMPLOYEE} & 1 & 9 & Low & 7 \\
\textbf{ILF-4} & \texttt{DOCUMENT} & 1 & 11 & Low & 7 \\
\textbf{ILF-5} & \texttt{LOAN} & 1 & 11 & Low & 7 \\
\textbf{ILF-6} & \texttt{REPAYMENT} & 1 & 12 & Low & 7 \\
\textbf{ILF-7} & \texttt{LOAN\_CATEGORY} & 1 & 7 & Low & 7 \\
\textbf{Total ILF Points} & - & - & - & - & \textbf{49} \\
\bottomrule
\end{tabular}%
}
\end{table}

\section{External Interface Files (EIF) Identification}

\subsection{Table 6.5: External Interface Files (EIF) Complexity Classification}

\begin{table}[htbp]
\centering
\small
\resizebox{\textwidth}{!}{%
\begin{tabular}{l l l l l l}
\toprule
\textbf{EIF ID} & \textbf{External System Interface File} & \textbf{Record Element Types (RET)} & \textbf{Data Element Types (DET)} & \textbf{Assessed Complexity} & \textbf{Weight} \\
\midrule
\textbf{EIF-1} & External Credit Bureau Score Interface (CIBIL/Experian) & 1 & 8 & Low & 5 \\
\textbf{EIF-2} & Payment Gateway Settlement Webhook Feed & 1 & 9 & Low & 5 \\
\textbf{Total EIF Points} & - & - & - & - & \textbf{10} \\
\bottomrule
\end{tabular}%
}
\end{table}

\section{Unadjusted Function Points (UFP) Calculation}

The Unadjusted Function Points (UFP) represents the aggregate sum of functional points across all five IFPUG categories:

\_\_DISPLAY\_MATH\_0\_\_

\_\_DISPLAY\_MATH\_0\_\_

\section{General System Characteristics (GSC)}

The IFPUG model defines 14 General System Characteristics (GSCs) representing global non-functional and environmental factors. Each characteristic is assigned a Degree of Influence (DI) rating from 0 (No Influence) to 5 (Strong / Essential Influence):

\subsection{Table 6.6: General System Characteristics Rating}

\begin{table}[htbp]
\centering
\small
\resizebox{\textwidth}{!}{%
\begin{tabular}{l l l l}
\toprule
\textbf{No.} & \textbf{Characteristic Description} & \textbf{Degree of Influence (0 - 5)} & \textbf{Justification} \\
\midrule
\textbf{C1} & Data Communications & 4 & Heavy online data exchange over HTTPS with distributed clients \\
\textbf{C2} & Distributed Data Processing & 3 & Distributed 3-tier architecture with separate application and DB tiers \\
\textbf{C3} & Performance Objectives & 4 & Strict transaction response constraints under 2 seconds for financial UI \\
\textbf{C4} & Heavily Used Configuration & 3 & High volume of concurrent loan applications and batch EMI updates \\
\textbf{C5} & Transaction Rate & 4 & High transaction frequency during banking hours and month-end EMI dates \\
\textbf{C6} & Online Data Entry & 5 & 100\% of customer proposals, documents, and reviews occur online \\
\textbf{C7} & End-User Efficiency & 4 & Optimized responsive workflows designed to accelerate underwriter tasks \\
\textbf{C8} & Online Update & 4 & Immediate real-time updates to loan balances, application states, and ledgers \\
\textbf{C9} & Complex Processing & 4 & Mathematical compound amortization formulas, risk algorithms, score logic \\
\textbf{C10} & Reusability & 3 & Modular components designed for multi-product expansion \\
\textbf{C11} & Installation Ease & 2 & Standardized automated web deployment; minimal client footprint \\
\textbf{C12} & Operational Ease & 4 & Automated daily backup procedures, audit logging, recovery mechanisms \\
\textbf{C13} & Multiple Sites & 3 & Multi-branch banking access across geographic regional offices \\
\textbf{C14} & Facilitate Change & 4 & Configurable loan categories, dynamic interest rates, flexible parameters \\
\textbf{Total Degree of Influence (TDI)} & - & \textbf{51} & Sum of all 14 General System Characteristics \\
\bottomrule
\end{tabular}%
}
\end{table}

\section{Value Adjustment Factor (VAF) and Final Function Points}

The Value Adjustment Factor (VAF) scales the unadjusted size based on the Total Degree of Influence:

\_\_INLINE\_MATH\_0\_\_	ext{VAF} = 0.65 + \textbackslash{}left(0.01 	imes \textbackslash{}sum\_{i=1}\textasciicircum{}{14} C\_i

ight) = 0.65 + (0.01 	imes 51) = 0.65 + 0.51 = \textbackslash{}mathbf{1.16}\_\_INLINE\_MATH\_0\_\_

The final Adjusted Function Points (FP) is determined as:

\_\_DISPLAY\_MATH\_0\_\_

\section{Effort and Schedule Estimation (COCOMO II Derivation)}

Using standard Software Engineering metrics (QSM / IFPUG benchmarks), one Function Point in an enterprise object-oriented language (Java / Python) corresponds to approximately 50 Source Lines of Code (SLOC):

\_\_DISPLAY\_MATH\_0\_\_

Applying the basic COCOMO (Constructive Cost Model) for Semi-Detached software projects:

\begin{itemize}
  \item \_\_INLINE\_MATH\_0\_\_
  \item \_\_INLINE\_MATH\_0\_\_
  \item \_\_INLINE\_MATH\_0\_\_
\end{itemize}

This mathematical estimation confirms that the Bank Loan Management System represents a robust, highly scoped engineering project ideally suited for academic and commercial enterprise development.
